\section{The Vision Beyond Nature and the Seam of Revelation}

Once created goods are weighed and found incapable of quieting the human will, the mind is tempted to settle for a dignified compromise: perhaps ultimate human fulfillment consists in philosophical contemplation---knowing \emph{that} God exists as the First Cause of the cosmos, tracing his wisdom through created effects, and living virtuously beneath his providence. That was the highest horizon available to classical pagan philosophy. St.\ Thomas Aquinas inspects that horizon in \emph{Summa Theologiae} I-II, q.\,3, a.\,8, and judges it still incomplete.

The reason lies in how the human intellect works. Every power achieves its perfection only when it reaches its proper object, and the proper object of the intellect is \emph{what a thing is}---its essence. Whenever the mind observes an effect and discovers that the effect has a cause, yet does not know what that cause is in itself, wonder remains alive and inquiry refuses to stop. Aquinas illustrates the point with a solar eclipse: a spectator who sees the sun darkened and realizes something must be causing the shadow wonders at the cause and searches until the nature of the eclipse is understood. The same law governs the mind's encounter with creation:

\begin{quote}
If therefore the human intellect, knowing the essence of some created effect, knows no more of God than ``that He is''; the perfection of that intellect does not yet reach simply the First Cause, but there remains in it the natural desire to seek the cause. Wherefore it is not yet perfectly happy. Consequently, for perfect happiness the intellect needs to reach the very Essence of the First Cause. And thus it will have its perfection through union with God as with that object, in which alone man's happiness consists, as stated above (AA.\ 1, 7; Q.\ 2, A.\ 8).\endnote{Aquinas, \emph{ST} I-II, q.\,3, a.\,8, corpus. In \emph{ST} I, q.\,12, a.\,1, corpus, Aquinas adds a theological and a rational ground for the possibility of this vision: denying that any created intellect can see God's essence either makes ultimate beatitude impossible or locates it outside God, contrary to faith; and it would leave the intellect's natural wonder at the First Cause perpetually unfulfilled.}
\end{quote}

Perfect beatitude, therefore, cannot rest in inferences drawn from finite mirrors. It requires seeing the divine essence itself. Yet the moment Aquinas names that requirement, he places a second truth directly beside it---one that shatters every illusion of human self-elevation.

\subsection*{Why No Creature Can See God by Nature}

In \emph{Summa Theologiae} I, q.\,12, aa.\,1, 4, and 5, Aquinas asks whether any created intellect---human or angelic---can see the essence of God by its own natural power. His answer is an unqualified no.

The argument turns on a fundamental principle of knowledge: the thing known exists in the knower according to the knower's own mode of being. Every creature has a mode of being distinct from its act of existence; even an angel, free from matter, is a subsisting form that \emph{receives} existence rather than \emph{being} existence itself. Our human intellect, bound in this life to a body, naturally abstracts universal natures from material things; an angelic intellect naturally grasps immaterial created natures. God alone is self-subsistent being (\emph{ipsum esse subsistens}), pure act with no admixture of potency. Therefore, to know subsistent being directly is natural to the divine intellect alone:

\begin{quote}
It follows therefore that to know self-subsistent being is natural to the divine intellect alone; and this is beyond the natural power of any created intellect; for no creature is its own existence, forasmuch as its existence is participated. Therefore the created intellect cannot see the essence of God, unless God by His grace unites Himself to the created intellect, as an object made intelligible to it.\endnote{Aquinas, \emph{ST} I, q.\,12, a.\,4, corpus. In reply 3, Aquinas notes why this elevation is possible for an intellect though impossible for a bodily sense: sense is wholly tied to singular matter, whereas intellect is naturally immaterial and capable of distinguishing form and existence abstractedly, so it possesses an openness to be raised above its native ceiling by grace.}
\end{quote}

Because God is pure act, he is in himself supremely knowable; the obstacle lies in the disproportion of our finite minds, like a bat blinded by noon sunlight (\emph{ST} I, q.\,12, a.\,1). Nor can any created image or concept bridge the gap, for no finite likeness can represent the infinite essence of God as he is in himself (\emph{ST} I, q.\,12, a.\,2; q.\,12, a.\,4, ad 1). When the blessed behold God, the divine essence itself becomes the intelligible form uniting with the created intellect.

To bear such union without being blinded or annihilated, the intellect must be strengthened by a supernatural disposition poured into it by God: what Catholic theology calls the \emph{lumen gloriae}, the light of glory (\emph{ST} I, q.\,12, a.\,5). Aquinas clarifies two vital boundaries about this created light. First, it is not a medium \emph{in which} God is seen, as though it were a mirror or screen standing between the soul and the divine face; it is a medium \emph{by which} the intellect is fortified to look directly upon God, preserving the immediacy of the vision (\emph{ST} I, q.\,12, a.\,5, ad 2). Second, by this light the blessed are made \emph{deiform}---God-like---according to 1 John 3:2, yet without ever becoming divine by nature or exhausting God's infinity. In \emph{ST} I, q.\,12, a.\,7, Aquinas distinguishes \emph{seeing} from \emph{comprehending}. To comprehend a reality strictly is to know it as completely as it is intrinsically knowable. Because God is infinite and every created light of glory remains finite, no creature can ever know God infinitely. The blessed see the \emph{whole} God---since God is utterly simple and has no hidden parts---yet they do not see him \emph{wholly} in the infinite degree proper to God's own self-knowledge.\endnote{Aquinas, \emph{ST} I, q.\,12, a.\,7, corpus and ad 1--3. In ad 1, Aquinas allows a broader sense of ``comprehension'' meaning secure attainment or possession (responding to hope, as vision responds to faith and enjoyment to charity), while strictly excluding any finite containment of God's infinity.}

\subsection*{The Seam Between Philosophy and Revelation}

Natural reason travels that far, and no farther. A philosophical ascent from created being and human desire can establish two weighty conclusions: first, that no finite good can quiet the rational will; second, that knowing God merely as First Cause from effects leaves the mind's wonder at \emph{what} God is still unfulfilled. Philosophy, however, cannot cross the threshold from the inadequacy of finite goods to the promise of the beatific vision. Still less can the natural desire to know the First Cause constitute a claim, right, or exigency that binds God to bestow supernatural union. A creature who owes its very existence to free creation can never present the Creator with an invoice for intimacy with his inner life.

To know that God has actually ordained human beings to see his face requires crossing a distinct seam---from natural reason to divine revelation. The First Vatican Council defined this boundary in the Dogmatic Constitution \emph{Dei Filius} (24 April 1870), Chapter 2, \emph{De Revelatione} (On Revelation). Having first affirmed that God, the beginning and end of all things, can be known with certainty by the natural light of human reason through created things (Rom 1:20), the Council teaches that divine revelation is \emph{absolutely necessary} precisely because God---out of his infinite goodness (\emph{ex infinita bonitate sua}), not out of any debt to human nature---has ordered humanity to an end above nature:

\begin{quote}
This, however, is not the reason why revelation is to be called absolutely necessary; but because God of His infinite goodness, has ordained man to a supernatural end, viz., to be a sharer of divine blessings which utterly exceed the intelligence of the human mind; for eye hath not seen, nor ear heard, neither hath it entered into the heart of man, what things God hath prepared for them that love Him.\endnote{First Vatican Council, Dogmatic Constitution \emph{Dei Filius}, ch.\ 2, \emph{De Revelatione}, 24 April 1870 (English translation attributed to Cardinal Henry Edward Manning in \emph{Dogmatic Canons and Decrees}, New York: Devin-Adair, 1912, p.\ 221; Latin text: ``Non hac tamen de causa revelatio absolute necessaria dicenda est, sed quia Deus ex infinita bonitate sua ordinavit hominem ad finem supernaturalem, ad participanda scilicet bona divina, quae humanae mentis intelligentiam omnino superant''). See also Canon II.3, anathematizing the claim that human nature cannot be raised by divine power to a knowledge and perfection surpassing nature.}
\end{quote}

Sacred Scripture announces this revealed promise in a steady crescendo from the Psalter to the New Testament. In the Psalms, the believer already looks beyond temporal deliverance to an encounter of faces: ``My heart hath said to thee: My face hath sought thee: thy face, O Lord, will I still seek'' (Ps 26:8 DR; Ps 27:8 Heb.).\endnote{Direct English Scripture quotations throughout follow the Challoner revision of the Douay-Rheims Bible (Project Gutenberg eBook 1581); scriptural verses embedded inside quoted translations of Aquinas, Augustine, Bernard, or other authors retain the wording of the quoted translation. Where Vulgate/Douay-Rheims psalm numbering diverges from Hebrew numbering, both numbers are noted on first citation.} St.\ Paul draws the sharp line between every earthly expectation and the revealed gift, quoting Isaias: ``That eye hath not seen, nor ear heard: neither hath it entered into the heart of man, what things God hath prepared for them that love him'' (1 Cor 2:9). In the same letter he contrasts our present pilgrimage with the coming unveiling: ``We see now through a glass in a dark manner: but then face to face. Now I know in part: but then I shall know even as I am known'' (1 Cor 13:12). St.\ John binds this sight directly to adoptive sonship and likeness to God: ``Behold what manner of charity the Father hath bestowed upon us, that we should be called and should be the sons of God. ... Dearly beloved, we are now the sons of God: and it hath not yet appeared what we shall be. We know that when he shall appear we shall be like to him: because we shall see him as he is'' (1 John 3:1--2).

The Catholic Church has received this revealed promise as defined dogma. In the Apostolic Constitution \emph{Benedictus Deus} (29 January 1336), Pope Benedict XII solemnly defined that the souls of all the saints and faithful who die in charity with nothing left to purify already behold the divine essence by an intuitive and face-to-face vision---with no creature intervening as an object of sight---even before the general resurrection of the body, and by that vision enjoy the divine essence and are truly blessed with eternal life and rest.\endnote{Pope Benedict XII, Apostolic Constitution \emph{Benedictus Deus}, 29 January 1336 (Denzinger--H\"unermann, nn.\ 1000--1002).} The \emph{Catechism of the Catholic Church} (nos.\ 1023--1029) expounds this teaching as the consummation of the entire Christian vocation: heaven is the ultimate end and fulfillment of the deepest human longings, the state of supreme and definitive happiness in communion of life and love with the Holy Trinity, the Virgin Mary, the angels, and all the blessed, where God gives himself directly to be seen while remaining an inexhaustible mystery.

\subsection*{Incapacity Without Humiliation}

At first glance, the doctrine that human beings cannot achieve their ultimate end by natural power sounds like a wound to human dignity. Every lower creature---an oak tree, a swallow, a hound---possesses within its own nature the equipment required to reach its proper flourishing. Why should the rational creature alone be born unable to reach its perfection by its own strength?

Aquinas confronts this objection head-on in \emph{ST} I-II, q.\,5, a.\,5, and his two replies turn the complaint inside out. First, he answers the comparison with irrational animals (ad 2): a nature capable of attaining the \emph{perfect} good with help from outside is nobler than a nature capable only of an \emph{imperfect} good on its own. A patient whose constitution can be restored to complete health through a physician's medicine is better off than someone who needs no medicine only because complete health is forever out of reach. Our inability to force open heaven by native muscle is not a defect below the beasts; it is the exact measure of a destiny pitched infinitely above them.

Second, in reply 1, Aquinas reveals what God placed in human hands in place of self-sufficiency:

\begin{quote}
Just as nature does not fail man in necessaries, although it has not provided him with weapons and clothing, as it provided other animals, because it gave him reason and hands, with which he is able to get these things for himself; so neither did it fail man in things necessary, although it gave him not the wherewithal to attain Happiness: since this it could not do. But it did give him free-will, with which he can turn to God, that He may make him happy. ``For what we do by means of our friends, is done, in a sense, by ourselves'' (\emph{Ethic.}\ iii, 3).\endnote{Aquinas, \emph{ST} I-II, q.\,5, a.\,5, ad 1; see also ad 2 on the nobility of needing divine assistance for the perfect good, and the corpus citing 1 Cor 2:9 and \emph{ST} I, q.\,12, a.\,4.}
\end{quote}

That quiet borrowing from Aristotle's \emph{Nicomachean Ethics}---``what we do by means of our friends, is done, in a sense, by ourselves''---changes the meaning of creaturely dependence. Human dependence is the receptivity of a person invited into friendship rather than the helplessness of an object pushed from outside. We cannot manufacture beatitude, because beatitude is another Person giving himself. Because the end of every human \emph{why} is communion with the living God, the path to it is friendship received as grace.
